\section{Gates for shared and projected labels}\label{sec:prices}

Suppose each of two groups of sites holds a common symbol. Part~I
constructed an equality filter for these symbols. In the many-petal
argument a group may learn only a projection of its symbol, for
example a few bits of the majority among the active satellites. We
will recover the equality test from such projections at a cost
independent of the original alphabet. To use projections as inputs to
further tests, we first specify how a gate acts on a sum of basis inputs.

\subsection{Controls and returned tensors}

Fix a finite set of sites \(S\), a spectral point and its log price \(h\)
as in Section~\ref{sec:words}. Put \(h(0)=-\infty\).
At each site \(v\) reserve a basis register with vectors \(e_{x_v}\),
and put \(e_x=\bigotimes_{v\in S}e_{x_v}\).
A \emph{gate} with domain \(\mathcal D\) and tensor output \(F(x)\)
retains these controls. On an input with arbitrary auxiliary registers
it acts by
\[
   Y=\sum_{x\in\mathcal D}e_x\otimes Y_x
   \quad\longmapsto\quad
   F[Y]=\sum_{x\in\mathcal D}e_x\otimes Y_x\otimes F(x).
\]
All tensor products are grouped by site. The gate has \emph{cost} at
most \(c\) if \(h(F[Y])\le h(Y)+c\) for every such input \(Y\).
A \emph{promise} specifies the subset of \(\mathcal D\) on which we
require a particular value of \(F(x)\). Cost is charged throughout
\(\mathcal D\). In particular two gates may agree and cancel on a
promise even when no site can check that promise.

A local map controlled by a basis register has cost zero. It can copy
a control, compute a function of it, or erase the register by sending
each basis vector to \(1\). Retaining the controls lets a second gate
use the first gate's output while carrying all the other registers as
auxiliaries. Costs then add. Every basis tuple in the support of the
new controls must lie in the second gate's domain, including on
unpromised inputs. A prior gate may enforce a domain condition by
sending every term outside the domain to zero.

A control common to all sites is \emph{public}. Given such a control,
each site may use it at zero additional cost. To sum over a new public
\(d\)-bit label we may supply \(G_S(2^d)\) at cost \(d\) and erase the
label after using it. For a transmission from one site \(u\) to
another \(v\), it suffices to supply \(G_{\{u,v\}}(2^d)\) and select
the sender's label at \(u\), at cost \(d\,w(\{u,v\})\).

We also add and subtract gates with common output spaces. The sum of
tensors is a restriction of their direct sum, so
\begin{equation}\label{eq:sum-prices}
              2^{h(\sum_j X_j)}\le \sum_j 2^{h(X_j)} .
\end{equation}
A sign can be applied at any one site. Thus gates of costs \(c_j\) can
be summed at cost at most \(\log_2\sum_j 2^{c_j}\). For example,
given local bits \(b_1,\ldots,b_n\), \(n\le r\), the exact count test
\begin{equation}\label{eq:count-test}
              \mathbf 1_{\{\sum_i b_i=t\}}=[z^t]\prod_{i=1}^n z^{b_i}
\end{equation}
costs at most \(\log_2(n+1)\). At each of \(n+1\) field elements the
product can be formed locally. Interpolation expresses the coefficient
as a fixed linear combination of these values.

Price bounds also pass to polynomial limits. Suppose \(X(t)\) has
degree at most \(d\) and \(h(X(t))\le c\) for \(t\ne0\).
Interpolate \(X(t)^{\otimes n}\) at \(dn+1\) distinct nonzero scalars
to recover \(X(0)^{\otimes n}\). Equation~\eqref{eq:sum-prices} gives
\(nh(X(0))\le nc+\log_2(dn+1)\). Division by \(n\) and passage to the
limit give \(h(X(0))\le c\).

If a gate returns a fixed nonzero tensor \(Q\) along with the prescribed
output on the promise, we call \(Q\) its \emph{offset}. A gross cost
bound \(c\) gives a \emph{net cost} bound \(c-h(Q)\).
To add branches with different offsets \(Q_i\), supply
\(\bigotimes_{j\ne i}Q_j\) in branch \(i\) and order the factors
identically at each site. The common offset is then
\(Q=\bigotimes_j Q_j\). If branch \(i\) had net cost \(c_i\), its
padded gross cost is at most \(h(Q)+c_i\). Subtraction or summation
therefore has net cost at most \(\log_2\sum_i 2^{c_i}\), even if some
branches give zero on a particular input.

\begin{lemma}[Diagonal comparison]\label{lem:diagonal}
Let \(P_1,\ldots,P_r\) be disjoint nonempty site groups with
\(K=\bigcup_iP_i\), \(w(K)=1\), and \(s=1-\sum_i w(P_i)>0\).
Let \(\mathcal W\) be a sunflower-free family of \(N\ge1\) distinct
words of length \(k\). For each coordinate \(j\) assume a gate of net
cost at most \(f_j\), with one symbol replicated throughout each
\(P_i\) as its only input controls. Its domain includes the column of
every ordered \(r\)-tuple from \(\mathcal W\). On a constant column it
returns a fixed nonzero tensor \(Q_j\), and on a column with a repeat
and at least two distinct symbols it returns zero. Then
\[
                     s\log_2 N\le \sum_{j=1}^k f_j .
\]
\end{lemma}

\begin{proof}
Give each group an independent word index by supplying
\(\bigotimes_i G_{P_i}(N)\). Keeping the indices, compute the symbol
registers locally, apply the \(k\) gates and erase the symbol
registers. A tuple of indices that is not constant has a violating
column by sunflower-freeness and contributes zero, regardless of the gate
outputs on its other columns. Each constant tuple contributes
\(\bigotimes_j Q_j\). Thus the output is
\(G_K(N)\otimes\bigotimes_j Q_j\). Comparing its price with the initial
price \((1-s)\log_2 N\) and the gross cost
\(\sum_j(h(Q_j)+f_j)\) proves the bound.
\end{proof}

\subsection{Broadcasting and equality}

There are two ways a gate can receive a shared label. Sometimes its
domain guarantees that all source registers agree. Sometimes a
previous test furnishes agreement only on a promise. The first case
admits an exact rate bound. The second needs the controlled
constructions of Part~I, with their returned tensors.

\begin{lemma}[Broadcasting and equality]\label{lem:broadcast}
Fix a spectral point on \(S\) over an algebraically closed field,
disjoint nonempty \(P,Z\subseteq S\), and \(K=P\cup Z\).
For \(b\)-bit labels, broadcasting appends \(E_Z(x)\) when all sources
in \(P\) have label \(x\). Equality testing multiplies by
\(\mathbf 1_{\{x=y\}}\) when \(P\) has common label \(x\) and \(Z\)
has common label \(y\). Set
\[
 c_{\rm b}=w(K)-w(P),\qquad
 c_{\rm e}=\max\{0,w(K)-w(P)-w(Z)\}.
\]
Both tasks have cost zero at \(b=0\). For \(b\ge1\) the bounds are
as follows.
\begin{enumerate}
\item[(i)] If the indicated agreements hold throughout the domain,
the tasks have unit offset and costs at most \(c_{\rm b}b\) and
\(c_{\rm e}b\), respectively.
\item[(ii)] If the agreements hold on a promise, the tasks admit
gates with fixed nonzero offsets and net costs at most
\(c_{\rm b}b+O(\sqrt{b\log_2(2|K|)})\) and
\(c_{\rm e}b+O(\sqrt{b\log_2(2|K|)})\), respectively.
\item[(iii)] In (ii), if \(w(K)=w(P)\), broadcasting has net cost
\(O(\log_2(2|K|))\) for all \(b\).
\end{enumerate}
\end{lemma}

\begin{proof}
First consider the promised tasks. Lemma~\ref{lem:copy} gives the
broadcast in (ii), and Lemma~\ref{lem:segment} gives (iii).
For equality, use the filter in the proof of Theorem~\ref{thm:filter}.
If \(w(K)>0\), normalize the price on \(K\) by dividing by \(w(K)\).
For this normalization the proof bounds the contribution at each
length \(m\le b\) by
\(2^{1-g'm+C_0\sqrt{m\ell}}\), where
\(g'=(w(P)+w(Z)-w(K))/w(K)\) and \(\ell=\log_2(2|K|)\).
Taking the logarithm of the sum including the unit base case,
then multiplying by \(w(K)\le1\), gives the stated linear term
and an error \(O(\sqrt{b\ell}+\log_2(b+1))=O(\sqrt{b\ell})\).
If \(w(K)=0\), the direct equality filter on \(K\) has cost zero.
All these controlled restrictions retain their controls and act on
arbitrary auxiliary factors, so their gross price bounds hold on the
entire domain.

For (i), take any input \(Y\) supported on the agreed domain and write
\(Y'\) for its desired output with unit offset. In \(Y^{\otimes n}\),
concatenate the \(n\) labels at each site. On this input (ii) produces
\((Y')^{\otimes n}\otimes Q_n\) for a fixed nonzero offset \(Q_n\).
For equality this uses the fact that equality of the concatenations
is the product of the \(n\) equality predicates. With
\(c_*=c_{\rm b}\) or \(c_{\rm e}\), subtract \(h(Q_n)\) to get
\[
              n h(Y')\le n h(Y)+c_*nb+O(\sqrt{nb\ell}).
\]
Taking \(n\to\infty\) after division by \(n\) gives the claimed
gate bound for every \(Y\).
\end{proof}

\subsection{Equality from projected labels}

The majority symbol in a class of satellites may be unknown at every
designated output site. The voting procedure in Section~\ref{sec:vote}
will supply copies of any projection of that symbol. The following
lemma quantifies when these copies suffice to test equality of the
original symbols. The projection procedure may also return an
auxiliary tensor depending on the input. The equality test preserves
this tensor on success.

\begin{lemma}[Equality from projected labels]\label{lem:oracle}
Fix a control domain and a promise on which \(a\in\{2,3\}\) labels
\(y_1,\ldots,y_a\in\{0,1\}^b\) are specified. Let \(P_1,\ldots,P_a\)
be disjoint nonempty groups, \(K=\bigcup_iP_i\), and \(w(K)=1\).
For each \(0\le l\le b\), suppose a gate receives the original
controls and an arbitrary public function
\(\pi:\{0,1\}^b\to\{0,1\}^l\). Its cost on this extended domain is
at most \(j+\lambda l+t\), and its output on the promise is
\[
                     J\otimes\bigotimes_i E_{P_i}(\pi(y_i)).
\]
The format of \(J\) is fixed, and its value depends only on the
original controls and may be zero. Suppose \(j,t\ge0\),
\(0\le\lambda\le2\), and
\(\sum_i w(P_i)\ge1+\lambda+\sigma\) with \(0<\sigma\le1\).
For \(a=3\), assume also
\(w(P_i\cup P_{i'})\ge w(P_i)+w(P_{i'})\) for every pair \(i\ne i'\).
Then a gate of cost at most
\[
              h(Q)+j+O\!\left(t+\log_2(2|S|)
                           +\sigma^{-1}\log_2(2|K|)\right)
\]
returns \(Q\otimes J\,\mathbf 1_{\{y_1=\cdots=y_a\}}\) on the promise
for a fixed nonzero tensor \(Q\).
\end{lemma}

\begin{proof}
We construct a gate for equality of \(\pi(y_1),\ldots,\pi(y_a)\) for
every output length \(l\), with one offset \(Q_l\) independent of
\(\pi\) and of the original controls. At \(l=0\) the projected labels
are empty. The supplied gate returns \(J\) at cost \(j+t\), so take
\(Q_0=1\).

\smallskip\noindent\emph{Shorten the labels.}
For \(l\ge1\) write \(\pi(y_i)=(h_i,x_i)\) with one head bit \(h_i\)
and \(d=l-1\) tail bits \(x_i\). Supply two independent \(d\)-bit
masks \(g_0,g_1\) by full diagonals and retain their basis registers.
Every site computes the composed public projection
\[
                  y\longmapsto x\oplus g_h
                  \qquad\bigl((h,x)=\pi(y)\bigr).
\]
Apply the length-\(d\) equality gate to this projection and erase the
temporary function register. On agreement every pair of masks
succeeds, and the result is
\(J\otimes Q_d\otimes G_S(2^d)^{\otimes 2}\).
If all heads agree but two tails differ, every pair fails. The
remaining error consists of the mask pairs that make different heads
agree on one translated tail \(z\). Each such pair occupies its own
basis register, so we can subtract it there.

\smallskip\noindent\emph{Reproduce the erroneous masks.}
Sum over assignments of the head bits to the \(a\) groups in which
both bits occur. For a fixed assignment, use the supplied projection
gate at length \(l\) and check the assigned heads locally. Sum over a
public \(d\)-bit string \(z\). For each occupied head reconstruct
its mask as follows. A singleton group broadcasts its tail to \(K\).
For a head shared by two groups, apply the equality test of
Lemma~\ref{lem:broadcast}(ii) to their two tails. On equal tails,
broadcast their common value from the union of the two groups to
the remaining sites of \(K\), again using (ii). In either case the
resulting tail is shared on \(K\). Extend it to \(S\) using (iii),
since \(w(K)=w(S)=1\), and output the mask \(g_h=x_i\oplus z\)
at every site. A broadcast with no new targets can be omitted.

On the promise this produces a mask pair exactly when the translated
tails agree. The head assignment is unique and the given pair
determines one \(z\), so each erroneous mask occurs with coefficient
one, including in positive characteristic. All temporary labels are
erased. The tensor \(J\) is the same in this correction as in the
shortened call, since it depends only on the original controls.
Off the promise each communication gate still acts on every basis
tuple in its registers with the asserted cost.

\smallskip\noindent\emph{Charge the correction.}
Pad the two branches and the head assignments with the unused
communication offsets, and pad the correction also with \(Q_d\).
For \(Q_l\) take the product of \(Q_d\), all these communication
offsets and the two full mask diagonals. On agreement this tensor
is returned with \(J\). On disagreement subtraction cancels every
surviving basis pair. Thus the specified output is obtained on the
whole promise, with an offset independent of \(\pi\).

Set \(\ell=\log_2(2|K|)\). For one occupied head let \(I\) be its
set of groups. The net cost to test its tails and broadcast to \(S\)
is bounded by
\[
       \Big(1-\sum_{i\in I}w(P_i)\Big)d
          +O\!\left(\sqrt{l\ell}+\log_2(2|S|)\right).
\]
For two groups, the pair hypothesis makes the equality rate
\(w(\bigcup_{i\in I}P_i)-\sum_{i\in I}w(P_i)\). Broadcasting to \(K\)
adds \(1-w(\bigcup_{i\in I}P_i)\) per bit, giving the displayed
coefficient. For one group the whole coefficient is its broadcast
rate. Extension from \(K\) to \(S\) has the length-independent cost
in (iii).

There are at most six mixed head assignments. Each costs \(d\) for
\(z\) and \(j+\lambda l+t\) for the projection gate. The two masks
it reconstructs occupy registers whose full diagonals contribute
\(2d\) to \(h(Q_l)\). Thus, relative to \(j+h(Q_l)\), the total
correction costs at most
\begin{align*}
 & t+\lambda l+d+\Big(2-\sum_iw(P_i)\Big)d-2d
                   +O\!\left(\sqrt{l\ell}+\log_2(2|S|)\right)\\
 &\hspace{1em}\le
       t+O\!\left(\ell/\sigma+\log_2(2|S|)\right)-\sigma d/2 .
\end{align*}
Here \(\lambda(l-d)\le2\), and the width surplus absorbs the square
root by \(C\sqrt{(d+1)\ell}\le\sigma d/2+O(\ell/\sigma)\).
The shortened branch has the preceding length's excess over
\(j+h(Q_d)\). By Equation~\eqref{eq:sum-prices}, exponentiating the
bounds adds the correction at each successive length to the base value
\(2^t\). The sum of the factors \(2^{-\sigma d/2}\), \(d\ge0\), is
\(O(1/\sigma)\). Its logarithmic contribution is absorbed in
\(O(\ell/\sigma)\), giving the claimed uniform bound. Finally take
\(l=b\) and let \(\pi\) be the identity.
\end{proof}
